\chapter{Clifford transport at one constituent}
\label{ch:hyp}

The completed theory of \cref{ch:connection} needs an $\so3$ that acts on the spinor module $S=\R^8$, with an unbroken generator and a
vector on which it acts. In the framework the unbroken generator is given: it is the charge $\half J$, the rotation of the surplus pair. The
algebra is then the rotation algebra of a triple formed by the surplus pair and a third element $X$,
\begin{equation}
\label{eq:triple}
\mathrm{span}\bigl\{\half e_2e_3,\ \half e_3X,\ \half Xe_2\bigr\},\qquad H(n)=n_1X+n_2e_2+n_3e_3 ,
\end{equation}
where $X$ is a symmetric involution that anticommutes with $e_2$ and $e_3$, so that the algebra is compact. These conditions do not single
out one element: $e_1$ and the volume element $\Gamma$ of $\Cl(3,3)$ anticommute, so every $X_\alpha=\cos\alpha\,e_1+\sin\alpha\,\Gamma$
satisfies them. Two choices are distinguished. The clock $e_1$ gives the temporal rotations $\so3_+$ of \PI{Remark~9.63}(c), and it fails on
the spinors (\cref{thm:ownclock}). Up to sign, the volume element $\Gamma$ is the only choice that commutes with the interior Lorentz algebra
(\cref{prop:unique}), and the orientation selects $+\Gamma$. This chapter states the hypothesis that selects it and works out what it gives at one constituent.

\section{The hypothesis}
\label{sec:hyp}

The hypothesis is adopted, not derived, and it is tested by its consequences. We state it first informally.

\begin{hypothesis}[Clifford transport, informal statement]
\label{hyp:informal}
The Yang--Mills term belongs to $V^+$. In the quantum regime, where observation in the classical sense is obstructed by $\hbar$, the reading
of the four-plane as a length is replaced by its electromagnetic counterpart. This holds at every rung, conditional on the reading.
\end{hypothesis}

The hypothesis extends an assignment that \PI{Section~9.6} records: the $b$-field on the positive triple, $c$ on the negative triple, and the
mass in the temporal volume element $T=e_1e_2e_3$. The electromagnetic counterpart of the four-plane $E$ is its space of two-forms
$\Lambda^2E$, on which the pseudoscalar $\Omega$ acts by left multiplication as minus the Hodge star $G$ of the field reading
(\PI{Prop.~3.7}), with the orientation of \cref{sec:conv}. The precise form used throughout is the following; its items are labelled (A1) to
(A4).

\begin{definition}[Clifford transport, precise form]
\label{def:hyp}
Let $n\ge1$ and let $\Cl(V^+_n)$ be the Clifford algebra of the $3n$ temporal directions, generated by the clock $E_1$ and the set $P$.
\begin{enumerate}[label=(A\arabic*),leftmargin=2.6em,itemsep=2pt]
\item \emph{Transport.} Under the field reading of the four-plane, the whole of $\Cl(V^+_n)$ acts on the spinor module through the
homomorphism
\begin{equation}
\label{eq:sigma}
\sigma_n:\quad E_1\longmapsto\Gamma_n,\qquad X\longmapsto X\quad(X\in P):
\end{equation}
the clock is replaced by the volume element $\Gamma_n$ of $\Cl(3n,3n)$, which satisfies $\Gamma_n^2=\id$.
\item \emph{Yang--Mills.} On the temporal directions the Yang--Mills field is a connection with values in the image of the rotations,
\begin{equation}
\label{eq:gn}
\gn:=\mathrm{span}\bigl\{\half XY\ (X,Y\in P,\ X\ne Y),\ \ \half X\Gamma_n\ (X\in P)\bigr\}.
\end{equation}
\item \emph{Mass.} The generator of mass is the image of the temporal volume, $G_n:=\sigma_n(T_n)$, with a rate $\kappa$ where it is a
rotation. The mass of a rung is carried by its generator of mass, carried into $V^-$ by the volume $\omega_S$. Where this applies, it
prescribes a vacuum direction of the mass field of (A4); it does not determine the full vacuum. At the first rung $G_1\,\omega_S=-\gamma^0$ is the clock vector, which is symmetric and not a mass kernel; it
selects the clock direction, and its image under (A1), the Higgs vector $\kappa\,\sigma_1(E_1)=\kappa\Gamma$, gives the mass kernel
$\kappa\,\Gamma\omega_S=\kappa\,T$. At the second rung $G_2\,\omega_S$ is itself an antisymmetric mass kernel, and (A3) prescribes the bidoublet value
$\kappa\,G_2\,\omega_S$; the Majorana value $v_R\Phi_0$ of \cref{ch:ew} is specified independently (\cref{thm:psstage}), and the image of the
clock vector has no vacuum value there. At the third rung $G_3$ is an involution, with no rate,
and (A3) selects no vacuum.
\item \emph{The spatial directions and the mass.} No Yang--Mills connection has values on the spatial triple $V^-$ of the four-plane: the
Yang--Mills value in $V^-$ is the scalar $m$, and the mass is a Lorentz-scalar field that can transform nontrivially under an internal
gauge group once projected on the four-plane. In formulas, the mass is a field $m(x)$ with values in
\begin{equation}
\label{eq:Mspace}
\mathcal M_n:=\bigl\{M:\ M^\top=-M,\ \ \Lambda^\top M\Lambda=M\ \text{for all}\ \Lambda\in\mathrm{Spin}(E)\bigr\},
\end{equation}
with $\mathrm{Spin}(E)$ the spin group of the four-plane, so that $\half\psi^\top m\,\psi$ is a Lorentz scalar; under the internal gauge
group the mass transforms as $m\mapsto UmU^\top$. The spatial
directions of the other constituents, the set $\cN$, carry a Yang--Mills connection with values in the antisymmetric part of their Clifford
algebra,
\begin{equation}
\label{eq:sn}
\mathfrak s_n:=\bigl\{C\in\Cl(\cN):\ C^\top=-C\bigr\},
\end{equation}
which is $0$ at one constituent, $\su2_L\oplus\su2_R$ at the second rung and $\so8$ at the third. The spatial directions also carry $c$, the
conversion of scales of \PI{Section~9.6}, and with it observation in the classical sense.
\end{enumerate}
Under the displacement reading nothing is replaced: the rotations act through the algebra $\kn$ spanned by the $\half XY$ with $X\ne Y$
and the $\half XE_1$, for $X,Y\in P$, and the volume is $T_n$.
\end{definition}

The definition does not say that one of the two readings is the true one: it attaches $\kn$ to the displacement reading and $\gn$ to the
field reading. The two algebras are conjugate at every rung, and they differ in one respect, which is the form that the exclusion of
\PI{Prop.~3.6} takes on the spinors: $\gn$ commutes with the interior Lorentz algebra and $\kn$ does not (\cref{thm:everyrung}). At the
second rung $\su2_L$ is the half of $\mathfrak s_2$ that acts on the left-handed $\mathbf4$ (\cref{thm:rung2}); $\mathfrak s_2$ is the weak
interaction with its right-handed partner. The Lorentz action is throughout the spectator action of \PI{Section~5.2}, in which the four-plane
and $\lint$ are those of the first constituent, as at the second rung in \PI{Section~8.5.3}.

\begin{remark}[Where $\hbar$ enters]
\label{rem:hbar}
The magnetic part of a frame field, on $V^+$, and the electric part, on $V^-$, are conjugate in Maxwell's theory: the bracket of
\PI{Thm~3.15}(iv) pairs them through the curl acting on test fields. After quantisation the equal-time commutator of smeared electric and
magnetic fields is a distribution proportional to $\hbar$, nonzero for suitable components and smearing functions. By \PI{Thm~3.15}(i) a
constant that converts a field strength into a length contains $\hbar$. These facts motivate the phrase ``obstructed by $\hbar$''. They do
not derive the transport rule, which remains a hypothesis.
\end{remark}

\section{Two algebras}

\begin{proposition}
\label{prop:twoalg}
$\so3_+=\mathrm{span}\{\half e_2e_3,\half e_3e_1,\half e_1e_2\}$ and $\mathfrak g_1=\mathrm{span}\{\half J,\half e_3\Gamma,\half\Gamma e_2\}
=\mathrm{span}\{\half J,\half e_2\Omega,\half e_3\Omega\}$, obtained from it by $e_1\to\Gamma$, consist of antisymmetric matrices, close as
$\su2$, share $\half J$ and have the Casimir $-\tfrac34$ on $S$.
\checks{ax-1, gm-3}
\end{proposition}

\section{The constituent's own clock}

\begin{theorem}
\label{thm:ownclock}
\begin{enumerate}[label=(\roman*),leftmargin=2.2em,itemsep=1pt]
\item Of $\so3_+$ only $J$ commutes with the interior Lorentz algebra, and only $J$ keeps the four kinetic forms $K^\mu$; all of $\so3_+$
keeps the mass form $T$.
\item $\so3_+$ and $\lint$ generate an algebra of dimension $15$, all of $\so{3,3}$.
\item The commutators of the two clock rotations with the boosts span the six mixed bivectors $e_ue_{a+3}$, $u=2,3$. A covariant coupling
therefore needs these as generators: with the two clock rotations, a charged four-vector of generators.
\item If the clock is a field, the Dirac data are covariant: for a rotor $U$ of $\mathrm{Spin}(V^+)$ with rotation $R$,
$U(n\cdot e)U^{-1}=(Rn)\cdot e$, $U^\top K^a(Rn)U=K^a(n)$ and $U^\top TU=T$. The breaking $Y^\top K^a+K^aY=2\,Y(\gamma^0)\gamma^a$ of
\PI{Prop.~9.53}(i) is exactly the motion of the clock.
\end{enumerate}
\checks{ax-2, ax-3, ax-4, ax-5, ax-6, ax-7, ax-8}
\end{theorem}

The massive pair of $\so3_+$ therefore cannot be coupled to the spinors with the Lorentz algebra held fixed. The invariant form on the coset
generators of (iii) is indefinite (\PI{Prop.~12.10}), and once the clock is a field the boosts themselves depend on it; this is the caveat of
\PI{Prop.~12.11}, that the boosts of the selected spacetime are not symmetries of that sector, seen on the spinors. In the terms of
\cref{def:hyp}, $\so3_+=\mathfrak k_1$ is the algebra of the displacement reading, and it does not coexist with a fixed interior Lorentz
algebra.

\section{The transported algebra is the commutant}

\begin{theorem}[The commutant of the Lorentz algebra]
\label{thm:clifford1}
\begin{enumerate}[label=(\roman*),leftmargin=2.2em,itemsep=1pt]
\item $\Gamma=\Omega J$, $e_1=-TJ$ and $J\gamma^0=T$. Hence two rules coincide: replacing $T$ by $-\Omega$ with $J$ unchanged, and
replacing the clock $e_1$ by $\Gamma$. In particular the description of the clock in \PI{Section~7.1.2}, as the quotient $\gamma^0=-TJ$ of
the rest rotation by the charge rotation, becomes $\Omega J=\Gamma$.
\item $(\Gamma,e_2,e_3)$ are anticommuting symmetric involutions that commute with $\lint$. So $\sigma:e_1\mapsto\Gamma$, $e_2\mapsto e_2$,
$e_3\mapsto e_3$ extends to a homomorphism of $\Cl(V^+)$ into the commutant of $\lint$ in $M_8(\R)$.
\item The commutant has dimension eight, and $\sigma$ is an isomorphism of $\Cl(V^+)$ onto it. Grade by grade (\cref{fig:transport}), it
maps the clock vectors $n\cdot e$ to the symmetric elements $H(n)=n_1\Gamma+n_2e_2+n_3e_3$, the rotations $\so3_+$ to the antisymmetric
algebra $\mathfrak g_1=\half\,\mathrm{span}\{J,e_2\Omega,e_3\Omega\}$, and the volume $T$ to $-\Omega$. The antisymmetric part of the
commutant is $\uu1_\Omega\oplus\su2$, the algebra $\mathfrak g_{\rm massless}$ of \PI{Section~7.1}.
\end{enumerate}
\checks{gm-1, gm-2, gm-3, gm-4}
\end{theorem}

\begin{proof}
(i) Moving the pair $e_2e_3$ through $e_4e_5e_6$ costs no sign, so $\Omega J=e_1e_2e_3e_4e_5e_6=\Gamma$; $TJ=e_1(e_2e_3)^2=-e_1$; and
$J\gamma^0=e_2e_3e_1=T$. Substituting $T\to-\Omega$ in $e_1=-TJ$ gives $\Omega J=\Gamma$.
(ii) $\Gamma$ is the volume element of an algebra with six generators, so it anticommutes with each $e_a$; $\Gamma^2=\id$ and
$\Gamma^\top=\Gamma$ follow from $e_a^2=\pm\id$ and the signs of the transposes. Each of $\Gamma$, $e_2$, $e_3$ anticommutes with $e_1$,
$e_4$, $e_5$, $e_6$, hence commutes with their bivectors, which span $\lint$. The relations of $\Cl(3,0)$ hold for $(\Gamma,e_2,e_3)$, so
$\sigma$ extends to a homomorphism.
(iii) The associative algebra generated by $\lint$ is the even Clifford algebra of the four-plane, $\Cl^0(1,3)\cong\Cl(3,0)\cong M_2(\C)$, of
real dimension eight; it acts faithfully, since it is simple. As a module over it $\R^8$ is the sum of two copies of the simple module
$\C^2$, so its commutant is $M_2(\C)$, of real dimension eight. The algebra $\Cl(3,0)\cong M_2(\C)$ is simple, so $\sigma$ is injective, and
its image, of dimension eight, is the whole commutant. The grades follow from (i) and the definitions; the symmetry of each image is a
direct computation.
\end{proof}

In the field reading, then, what belongs to $V^+$ is exactly what commutes with the Lorentz algebra of the four-plane: its clock vectors,
its rotations and its volume, and nothing else. In particular, inside this commutant the third element of \cref{eq:triple} is fixed up to sign, and the transport, with the chosen
orientation, selects $\Gamma$.

\begin{proposition}[The third element]
\label{prop:unique}
The symmetric involutions in the commutant of $\lint$ that anticommute with $e_2$ and $e_3$ are $\pm\Gamma$. The sign is fixed by the image
of the oriented volume: with $+\Gamma$, $\sigma(T)=-\Omega$, which acts on the field strength as $G$. Of the elements
$X_\alpha=\cos\alpha\,e_1+\sin\alpha\,\Gamma$, which all satisfy the conditions of \cref{eq:triple}, only $\pm\Gamma$ commute with $\lint$.
\checks{ax-24}
\end{proposition}

\begin{proof}
By \cref{thm:clifford1}(iii) the symmetric elements of the commutant are the images of grades zero and one, $a\,\id+b\,\Gamma+c\,e_2+d\,e_3$.
Anticommuting with $e_2$ requires $a=c=0$, and with $e_3$ also $d=0$. An involution $b\,\Gamma$ has $b=\pm1$. Finally $e_1$ does not commute
with the boosts, so $X_\alpha$ commutes with $\lint$ only for $\cos\alpha=0$.
\end{proof}

\begin{figure}[t]
\centering
\begin{tikzpicture}[>=Stealth,font=\small]
\matrix (m) [matrix of nodes, row sep=8pt, column sep=22pt, nodes={anchor=center,inner sep=2pt},
  column 1/.style={nodes={font=\footnotesize,text=black!60}}] {
 grade & $\Cl(V^+)$ & commutant of $\lint$ & kernel $C\,\omega_S$\\
 $0$ & $1$ & $\id$ & $\omega_S$, symmetric\\
 $1$ & $n\cdot e$, clock vectors & $H(n)=n_1\Gamma+n_2e_2+n_3e_3$ & $H(n)\,\omega_S$, Higgs-valued\\
 $2$ & $\so3_+$, rotations & $\mathfrak g_1=\tfrac12\langle J,e_2\Omega,e_3\Omega\rangle$ & $Y\omega_S$, Yang--Mills-valued\\
 $3$ & $T$, volume & $-\Omega$, generator of mass & $-\gamma^0$, symmetric\\
};
\draw[->] (m-2-2) -- node[above,font=\footnotesize]{$\sigma$: $e_1\mapsto\Gamma$} (m-2-3);
\draw[->] (m-2-3) -- node[above,font=\footnotesize]{$\cdot\,\omega_S$} (m-2-4);
\foreach \r in {3,4,5} {\draw[->] (m-\r-2) -- (m-\r-3); \draw[->] (m-\r-3) -- (m-\r-4);}
\node[draw,dashed,rounded corners=3pt,fit=(m-3-4)(m-4-4),inner sep=3pt,label={[font=\footnotesize]right:$\mathcal M_1$}] {};
\end{tikzpicture}
\caption{Clifford transport at one constituent (\cref{thm:clifford1,thm:massfield}). The map $\sigma$ is an isomorphism of $\Cl(V^+)$ onto
the commutant of the interior Lorentz algebra; multiplication by the volume $\omega_S$ of the spatial triple turns the commutant into the
Lorentz-invariant kernels, of which the antisymmetric ones, of grades one and two, form the mass field $\mathcal M_1$. The image $-\Omega$
of the temporal volume acts on the field strength, by left multiplication, as the binding root (\cref{thm:massgen}).}
\label{fig:transport}
\end{figure}

\begin{theorem}[The same rule as a conjugation]
\label{thm:conjugation}
\begin{enumerate}[label=(\roman*),leftmargin=2.2em,itemsep=1pt]
\item $\Omega^2=-\id$, $\Omega$ commutes with $\lint$, and on the six bivectors of the four-plane, the rotations $\rho_i=\half\gamma^j\gamma^k$
and the boosts $\beta_i=\half\gamma^0\gamma^i$, $\Omega\rho_i=-\beta_i$ and $\Omega\beta_i=\rho_i$. By \PI{Prop.~3.7} this is minus the Hodge
star $G$ of the field reading.
\item $\mathfrak g_1$ commutes with $\lint$ and keeps the four kinetic forms; with $\lint$ it generates an algebra of dimension $9=6+3$. In
the matrix realisation of \PI{Prop.~9.54} it is the algebra historically associated with Pauli and G\"ursey \cite{Pauli57,Gursey58}.
\item Let $g=P_++JP_-$, $P_\pm=\half(\id\pm\omega_S)$, the map of \PI{Prop.~9.54}(iv). Then $g$ is orthogonal, $g\,e_1\,g^{-1}=\Gamma$,
$g\,e_2\,g^{-1}=-e_3$, $g\,e_3\,g^{-1}=e_2$ and $g\,T\,g^{-1}=-\Omega$: conjugation by $g$ is $\sigma$ followed by the quarter turn
$e_2\mapsto-e_3$, $e_3\mapsto e_2$ of the surplus pair, also on the rotations. For instance $g(\half e_3e_1)g^{-1}=\half e_2\Gamma$, while
$\sigma(\half e_3e_1)=\half e_3\Gamma$. The images of the two algebras agree, $g\,\so3_+g^{-1}=\mathfrak g_1$, and the orthogonal map
$e^{-\pi J/4}g$ implements $\sigma$ exactly.
\item $\mathfrak g_1$ is the rotation algebra of $(\Gamma,e_2,e_3)$: in \cref{eq:triple}, $X=\Gamma$.
\end{enumerate}
\checks{ax-2, ax-3, ax-4, ax-13, ax-14, ax-15, ax-22, ax-25}
\end{theorem}

Under the stated choices the transported algebra, the algebra of Pauli and G\"ursey and the rotation algebra of the triple with $X=\Gamma$
therefore coincide at one constituent, as image algebras. The algebra does not fix the sign: $\Gamma$ and $-\Gamma$ generate the same rotation
subalgebra, and the orientation convention selects $+\Gamma$. Nor does it imply the dynamical assignments (A3) and (A4).

\section{The generator of mass}

\PI{Section~7.1.2} describes two rotations. Dirac matter at rest moves by $\Psi(\lambda)=e^{-m\lambda T}\Psi_0$, the one-parameter group of
the temporal volume element, and the frame moves by the binding rotation $e^{-\kappa G\tau}$. Part~I finds that the two do not meet: on frame
vectors $T$ acts as parity (\PI{Prop.~7.5}), $m$ and $\kappa$ are independent inputs, and reading $\hbar\kappa$ as $mc^2$ ``would be an
added identification, not a result''.

\begin{theorem}[The rest rotation becomes the binding rotation]
\label{thm:massgen}
Write the field strength of the four-plane as $F_X=\sum_i(B_i\rho_i+E_i\beta_i)$, the magnetic part on the rotations $\rho_i$ and the
electric part on the boosts $\beta_i$ of \cref{thm:conjugation}.
\begin{enumerate}[label=(\roman*),leftmargin=2.2em,itemsep=1pt]
\item Left multiplication by $\sigma(T)=-\Omega$ is the binding root: $-\Omega\,F_X=F_{GX}$.
\item $e^{\theta\Omega}F_X=F_{e^{-\theta G}X}$, again by left multiplication; conjugation by $e^{\theta\Omega}$ leaves $F_X$ unchanged, since
$\Omega$ commutes with the bivectors of the four-plane. So $\sigma$ carries the Dirac rest rotation $e^{-m\lambda T}$ to $e^{m\lambda\Omega}$,
which acting on the field strength is the binding rotation $e^{-\kappa G\tau}$ with $\kappa\tau=m\lambda$.
\item In the displacement reading the rotations of $V^+$ commute with $T$. In the field reading their images commute with $-\Omega$ and with
$\lint$, and of them only $J$ commutes with $T$.
\end{enumerate}
\checks{gm-5, gm-6, gm-10}
\end{theorem}

Under the hypothesis the two rotations are images of one element of $\Cl(V^+)$: $\sigma$ carries the generator $T$ of the Dirac rest rotation
to the generator $-\Omega$ of the binding rotation, and the angles correspond, $\kappa\tau=m\lambda$. This is a correspondence of generators,
not an equality of rates, and the matching of angles is kinematic. That the rest frequency of a fermion equals the rate of the binding
rotation, $\hbar\kappa=mc^2$, needs in addition that $\lambda$ and $\tau$ are the same physical time and that the normalisations of the
fields agree; the coupled action of \cref{sec:coupled} shows where the normalisations enter. The algebra thus reduces the identification of
Part~I to a normalisation, without deriving it. In the sense of (A3), at one constituent the generator of mass is $G_1=\sigma(T)=-\Omega$,
which acts on the field strength as $G$, and the mass is its rate $\kappa$.

\begin{proposition}[Two roles of $T$ come apart]
\label{prop:twojobs}
In the displacement reading $T$ has two roles: it generates the rest rotation, and it is the kernel of the mass term, invariant under the
rotations of $V^+$. In the field reading the two separate.
\begin{enumerate}[label=(\roman*),leftmargin=2.2em,itemsep=1pt]
\item The generator $-\Omega$ is invariant under $\mathfrak g_1$ and $\lint$.
\item The kernel space of \PI{Lemma~7.11} is $\gamma^0$ times the span of $J$, $e_2$, $e_3$, $e_2\Omega$, $e_3\Omega$ and $\Gamma$, six of the
eight elements of the commutant. No nonzero kernel in it is invariant under $\mathfrak g_1$: every mass kernel breaks the Yang--Mills
algebra.
\item The rotation generated by $\Omega$ preserves the kinetic forms and sends the kernel $T$ to $Te^{2\xi\Omega}$ (\PI{Prop.~7.2}(iii)): the
kernel has charge two under the generator of mass.
\end{enumerate}
\checks{gm-6, gm-7, gm-9}
\end{proposition}

In the field reading the name ``generator of mass'' therefore belongs to the rotation generator $-\Omega$, which acts on the field strength
as $G$, together with its rate; the kernel is a second object, the value of the mass field of \cref{sec:massfield}, a vector of
$\mathfrak g_1$.

\section{The Higgs vector is the clock}

\begin{proposition}
\label{prop:higgs1}
Let $H(n)=\sigma(n\cdot e)=n_1\Gamma+n_2e_2+n_3e_3$.
\begin{enumerate}[label=(\roman*),leftmargin=2.2em,itemsep=1pt]
\item $H(n)^2=|n|^2$, $H(n)$ is symmetric and commutes with $\lint$, and for unit $n$, $\half(\id+H(n))$ is a projector of rank four. $H$
transforms as a vector under the group of $\mathfrak g_1$, and $J$ fixes only the poles $n=\pm(1,0,0)$, the projectors
$\half(\id\pm\Gamma)$.
\item The mass kernel $M(n)=H(n)\,\omega_S$ is antisymmetric and Lorentz invariant, and $\Gamma\omega_S=T$. It spans the triple
$(e_2\omega_S,e_3\omega_S,T)$ of the kernel space, which $\mathfrak g_1$ maps into itself; $\Omega$ maps it onto the other triple
$(e_1e_2,e_1e_3,J\omega_S)$.
\item The variations of $T$ under $\half J$, $\half e_3\Omega$, $\half e_2\Omega$ and $\half\Omega$ have squared norms $0,8,8,8$.
\item With the rate as length, $\sigma(\kappa e_1)=\kappa\Gamma$ has square $\kappa^2$, and its kernel $\kappa\Gamma\omega_S$ is $\kappa T$, the
Dirac mass kernel with $m=\kappa$ when the Yukawa coupling is one.
\end{enumerate}
\checks{ax-9, ax-10, ax-11, ax-16, ax-17, ax-18, gm-8}
\end{proposition}

By (A3) the Higgs vector of \cref{ch:connection} is the image of the clock vector, and its length is identified with the rate. The vacuum
direction is then not chosen: it is the clock, $H_0=\kappa\Gamma$. How the length becomes masses depends on normalisations that the algebra
does not fix; they are kept in the coupled action of \cref{sec:coupled}. The kernel field of \PI{Chapter~7} and the projector field of
\PI{Chapter~8} are one vector, related by $M=H\omega_S$. The projector has fixed length, and by \cref{thm:bounded} that is not enough: the
completion needs a length that is a field. What fixes the value $\kappa$ is open (\PI{Prop.~4.49}).

\begin{proposition}[The massless criterion at the first rung]
\label{prop:massless1}
At the pole $\Pi_0=\half(\id+\Gamma)$, with $Q=\half J$: $[Q,\Pi_0]=0$; the two tangents $\delta\Pi=[\half e_3\Omega,\Pi_0]$ and
$[\half\Omega e_2,\Pi_0]$ are symmetric, charged and orthogonal, with $\Tr\delta\Pi=0$ and $\Tr[(\delta\Pi)^2]=2$ for each; and the mass
pattern $\Tr([Y,\Pi_0]^2)$ along the axes $(e_1,e_2,e_3)$ is $(0,2,2)$. The two tangents form a real charge plane, of which complex
combinations are the charge eigenstates. No bivector of $\so{3,3}$ moves $\Pi_0$.
\checks{ax-19, ax-20, ax-21}
\end{proposition}

The component of the connection along the clock axis is massless and the two along the surplus pair are massive: the surplus pair gives the
two directions into which the Higgs vector tilts, and these are the modes absorbed by the massive pair. The criterion of \PI{Cor.~8.53} is
met at the first rung by a Lorentz-invariant vacuum once variations by $\half e_2\Omega$ and $\half e_3\Omega$ are admitted;
\PI{Prop.~8.55} concerns variations generated by bivectors and stands.

\section{The mass field}
\label{sec:massfield}

By (A4) the mass is a Lorentz-scalar field that can transform nontrivially under an internal gauge group once projected on the four-plane.
At one constituent this has an exact form.

\begin{theorem}[The mass as a Lorentz-scalar field]
\label{thm:massfield}
\begin{enumerate}[label=(\roman*),leftmargin=2.2em,itemsep=1pt]
\item Under all of $\so{3,3}$ the invariant antisymmetric kernels are the multiples of $T$ (\PI{Cor.~7.10}), and the invariant symmetric
ones form a line. Under the Lorentz algebra of the four-plane alone, the invariant antisymmetric kernels form the six-dimensional space
$\mathcal M_1$ of \cref{eq:Mspace}, and the symmetric ones a plane. Projected on the four-plane, the single mass $T$ becomes one point of
$\mathcal M_1$.
\item Every Lorentz-invariant kernel is $C\omega_S$ with $C$ in the commutant of $\lint$, so $C=\sigma(c)$ with $c\in\Cl(V^+)$. The kernels
$\sigma(c)\,\omega_S$ are symmetric for $c$ of grades zero and three, $\omega_S$ and $\sigma(T)\,\omega_S=-\gamma^0$, and antisymmetric for
grades one and two. Hence $\mathcal M_1=\sigma(V^+)\,\omega_S\oplus\mathfrak g_1\,\omega_S$: the Higgs vectors and the Yang--Mills generators,
each multiplied by the volume of $V^-$ (\cref{fig:transport}).
\item $\mathfrak g_1$ acts on $\mathcal M_1$ as on $\Cl(V^+)$, $[Y,\sigma(c)\omega_S]=\sigma([y,c])\,\omega_S$. The Higgs-valued triple is the
vector and the Yang--Mills-valued triple the adjoint of $\su2$; no nonzero element of $\mathcal M_1$ is invariant.
\item The generator of mass acts on $\mathcal M_1$ with charge two, $M\mapsto Me^{2\xi\Omega}$, and $\Omega$ maps the Higgs-valued triple onto
the Yang--Mills-valued one. With $\Omega$ as the imaginary unit, $\mathcal M_1$ is the complex vector of $\su2$, the module of
\PI{Thm~7.12}.
\item The elements of $\mathcal M_1$ that commute with the charge $J$ are $T=\Gamma\omega_S$, the Higgs value at the pole, and
$T\Omega=J\omega_S$, the Yang--Mills value along the unbroken generator: the scalar and the pseudoscalar Dirac mass of \PI{Prop.~7.2}(ii).
The other four anticommute with $J$.
\end{enumerate}
\checks{mf-1, mf-2, mf-3, mf-4, mf-5}
\end{theorem}

\begin{proof}[Proof of the first sentence of (ii)]
For $X\in\lint$ one has $X^\top\omega_S=-\omega_SX$: the rotations are antisymmetric and commute with $\omega_S$, and the boosts are
symmetric and anticommute with it. Hence $X^\top(C\omega_S)+(C\omega_S)X=[X^\top,C]\,\omega_S$, which vanishes for all $X\in\lint$ exactly
when $C$ commutes with $\lint$, since $\lint$ is closed under transposition.
\end{proof}

The mass is thus a field $m(x)\in\mathcal M_1$. It is a Lorentz scalar, since $\half\psi^\top m\,\psi$ is invariant under the Lorentz
algebra of the four-plane, and it transforms nontrivially under the internal gauge algebra, as the complex vector of $\su2$. ``Projected on
the four-plane'' is the operative condition: under all of $\so{3,3}$ only $T$ is invariant, in the sense that the bivector action preserves
it as a bilinear form. Under the transported internal algebra, $T=\Gamma\omega_S$ belongs to the Higgs triplet, and its stabiliser in
$\mathfrak g_1$ is the line generated by $J$.

\begin{proposition}[The mass in every direction, and the component neutral under the charge]
\label{prop:massread}
For every $h\in\R^3$, $m(h):=H(h)\,\omega_S$ satisfies $m(h)^\top m(h)=|h|^2$ and anticommutes with $\alpha^a=\gamma^0\gamma^a$, $a=1,2,3$,
which satisfy $(\alpha^a)^2=\id$ and anticommute with one another. Hence, for a constant $h$,
\[
\bigl(\alpha^a\partial_a+m(h)\bigr)^2=\Delta-|h|^2 :
\]
the fermion mass is $|h|$ in every direction of the Higgs vector. For a field $h(x)$,
\[
\bigl(\alpha^a\partial_a+m(h(x))\bigr)^2=\Delta-|h|^2+\alpha^a\,\partial_am(h(x)) :
\]
the cross terms that differentiate the spinor cancel, and those that differentiate the mass field remain, so a dispersion law and a pole mass
refer to a constant background. The part of $m(h)$ that commutes with the charge $J$ is the clock component $h_1T$, a Dirac mass, and the
rest anticommutes with $J$. A rotation of the Higgs direction with the axis of $J$ held fixed keeps $|h|$ and changes the component neutral
with respect to that axis; the rotation with cosine $-\tfrac7{25}$ turns $T$ into $-\tfrac7{25}T\mp\tfrac{24}{25}e_2\omega_S$. Rotating the
Higgs vector together with its stabiliser is a gauge transformation and changes no spectrum.
\checks{mf-6, mf-7}
\end{proposition}

This is the spinor form of \PI{Prop.~9.58}, where under a temporal $b$-field the four-dimensional mass read by a fixed clock oscillates at a
fixed six-dimensional mass: here the length $|h|$, which gives the pole mass $|y|\,|h|$ on a constant real-triplet background, plays the
part of the six-dimensional mass, and the component $h_1T$ neutral under the fixed reference $J$ that of the four-dimensional one.

\section{The coupled action at one constituent}
\label{sec:coupled}

With every normalisation kept, the action of the spinor, the connection and the Higgs vector at one constituent is
\begin{equation}
\label{eq:coupled}
L=\half\Psi^\top K^\mu D_\mu\Psi-\tfrac y2\,\Psi^\top m(h)\,\Psi-\tfrac14F^A_{\mu\nu}F^{A\mu\nu}+\tfrac Z2\,D_\mu h\cdot D^\mu h
-\tfrac\lambda4\,(h\cdot h-v^2)^2 ,
\end{equation}
with $D_\mu\Psi=\partial_\mu\Psi+g\,b^A_\mu t_A\Psi$, where $t_A=\half(J,\ e_3\Gamma,\ \Gamma e_2)$ are the generators of $\mathfrak g_1$
about the axes $(e_1,e_2,e_3)$, $D_\mu h=\partial_\mu h-g\,b_\mu\times h$, and $m(h)=H(h)\,\omega_S$. The sign in $D_\mu h$ is fixed by
covariance, $H(D_\mu h)=\partial_\mu H(h)+g\,[b^A_\mu t_A,H(h)]$, since $[t_A,H(h)]=-H(\hat e_A\times h)$; it is the convention of
\cref{ch:connection} with $g$ replaced by $-g$, which changes no mass. $\Psi$ is a Grassmann field, as in \PI{Prop.~7.2}. The coefficients
$y$, $Z$, $\lambda$, $g$ and $v$ are independent.

The action restricts the mass field to the three-dimensional subspace $\sigma(V^+)\,\omega_S$ of Higgs-valued kernels, half of the
six-dimensional $\mathcal M_1$ of \cref{thm:massfield}; this is a choice of the model. The single-mass property of \cref{prop:massread} does
not extend to all of $\mathcal M_1$: the kernel $T+(e_2\omega_S)\,\Omega$ lies in $\mathcal M_1$, and its singular values are $0$ and $2$,
each four times.
\checks{mf-8}

\begin{proposition}[What the action gives]
\label{prop:coupled}
\begin{enumerate}[label=(\roman*),leftmargin=2.2em,itemsep=1pt]
\item $L$ is invariant under finite internal rotations, $\Psi\to U\Psi$ and $h\to Rh$ with $UH(h)U^{-1}=H(Rh)$, and under the Lorentz group of
the four-plane: $U^\top K^\mu U=K^\mu$, $U\,m(h)\,U^\top=m(Rh)$ and $\Lambda^\top m(h)\Lambda=m(h)$.
\item The field equation of $\Psi$ is $K^\mu D_\mu\Psi=y\,m(h)\Psi$. In the vacuum $h=(v,0,0)$, $b=0$, a fermion at rest moves by
$\Psi(t)=e^{-yvTt}\Psi_0$, a rotation generated by $T$ with the oriented rate $yv$; the physical mass is $m_f=|y|v$ for $v>0$. At every momentum
$\ddot\Psi=(\Delta-y^2v^2)\Psi$.
\item The component of $b$ along the clock axis is massless; the two along the surplus pair have $M_W^2=Zg^2v^2$; the radial mode of $h$ has
$M_h^2=2\lambda v^2/Z$.
\item The currents $\half\Psi^\top K^\mu t_A\Psi$ are nonzero Grassmann bilinears, and on solutions of the free fermion equation their
divergence is $y\,\Psi^\top t_A\,m(h)\Psi$, which vanishes when $t_A$ commutes with the mass. In the vacuum the current of the charge
$Q=\half J$ is conserved and the two broken currents are not. $Q$ has the charges $\pm\half$ on the spinors and $\pm1$ on the massive pair.
\item In unitary gauge, on a region where $h\ne0$ and $h=(v+\rho)\,\hat e_1$, the massless component is the clock component $a=b^1$ of the
connection, with curvature $f=da$, as in \cref{prop:whatadds} with the clock axis in the place of the third. The frame vector $X=(\bB;\bE)$
of $f$ obeys the system of \PI{Prop.~3.10} with no magnetic current, since $df=0$, and with the electric current of (iv) and the current of
the massive pair as source; $h$ has no current along its own axis. The clock component of the full curvature differs from $f$ by a term
proportional to $g\,b^2\wedge b^3$, quadratic in the massive pair, and obeys only the covariant identity of \cref{thm:frame}: its ordinary
Bianchi identity has a source made of the massive pair. The gauge-invariant form of $f$ is the tensor of 't~Hooft \cite{tHooft74}, which
reduces to $f$ in unitary gauge. In the terms of \cref{ch:curl}, $A_+X=-(0;\bj)$, so with the momentum convention there $\pi=-m(0;\bj)$. Its
Lagrangian is Maxwell's, $-\half\eta(X,X)$, the pure kinetic term of the residual connection, and not the split form of \cref{ch:curl}; the
expansion of \cref{prop:whatadds} also contains terms quadratic in $a$ and quadratic in $W$, from $|D\wedge W|^2$.
\end{enumerate}
\checks{ca-1, ca-2, ca-3, ca-4, mf-6, co-5, for the algebra of (i) to (iv) and the sign of $D_\mu h$; the field equations follow by the
variation of \PI{Prop.~7.2}, with \cref{thm:frame} and \cref{prop:whatadds}}
\end{proposition}

\paragraph{The two rates.} The rest evolution of a fermion is generated by $T$, the vacuum value of the mass field, in both readings. The
binding rotation of the field strength is generated by $-\Omega=\sigma(T)$. The vacuum kernel $T$ enters the Yukawa term of
\cref{eq:coupled}. The transported generator $-\Omega$ acts on field-strength bivectors and rotates the full kernel space $\mathcal M_1$
(\cref{thm:massfield}(iv)); a generic such rotation takes the Yukawa kernel outside the real triplet retained in this action, so it is not a
symmetry of \cref{eq:coupled}. The equality of the rates is the statement $|y|v=\kappa$: the rest frequency of the fermion equals the rate of
the binding rotation. By (A3), $v=\kappa$. We take in addition $y=1$, and in general $y\ge0$. Both
identifications refer to the normalisation of \cref{eq:coupled}: in the canonical scalar variable $h_c=\sqrt Z\,h$ the vacuum length is
$\sqrt Z\,v$ and the Yukawa coefficient is $y/\sqrt Z$. With them $m_f=\kappa$, $M_W=g\sqrt Z\,\kappa$ and $M_h=\sqrt{2\lambda/Z}\,\kappa$.
The correspondence of generators of \cref{thm:massgen} is exact; the equality $m_f=\kappa$ is conditional on (A3) and on $y=1$ in this
normalisation, and $g$, $Z$ and $\lambda$ are left to the dynamics.

\section{The Yang--Mills term and the doublet}

\begin{theorem}[The quadratic Lagrangians of the curvature]
\label{thm:ymforced}
The symmetric bilinear forms on a frame vector that are invariant under $\lG$ form a space of dimension two, spanned by $\eta$ and $\eta G$.
Hence a Lagrangian density that is quadratic in the curvature components $X^A$, invariant under the connected Lorentz group and invariant
under $\mathfrak g_1$, is
\[
a\sum_A\eta(X^A,X^A)+b\sum_A\eta(X^A,GX^A) :
\]
the Yang--Mills term of \cref{thm:frame} and the density $\bE^A\!\cdot\!\bB^A$, which is proportional to $\mathrm{tr}(F\wedge F)$.
\checks{ax-23, for the dimension}
\end{theorem}

The second density is locally the derivative of a Chern--Simons form, so it leaves the local classical equations in the bulk unchanged at
this order; its integral distinguishes topological sectors, and an independent angle multiplying it can affect the quantum theory. If parity
is imposed it is excluded. Given these invariances, the Yang--Mills term is fixed up to its scale as far as the local classical equations
are concerned.

\begin{proposition}[One doublet]
\label{prop:onedoublet}
The commutant of $\lint$ and $\mathfrak g_1$ on $S$ has dimension two, spanned by $\id$ and $\Omega$: $S$ is one irreducible module of
complex type, a single Weyl doublet.
\checks{ax-12, gm-4}
\end{proposition}

An $SU(2)$ gauge theory with an odd number of Weyl doublets is inconsistent \cite{Witten82}. This concerns the theory with the fermion
doublet; the bosonic model of \cref{ch:connection} has no such anomaly. At one constituent the hypothesis therefore gives the right algebra
on a module that cannot be gauged by itself as a quantum theory. Whether the higher rungs are consistent is decided by their own matter
content, which \cref{ch:rungs} examines (\cref{thm:ps}).

\section{Relation to Part~I}

\PI{Section~7.1.2} calls $\hbar\kappa=mc^2$ an added identification. Under the hypothesis the generators of the rest rotation and of the
binding rotation correspond, $\sigma(T)=-\Omega$ (\cref{thm:massgen}), and what remains to identify is a normalisation, $yv=\kappa$ in the
coupled action, which we assume; the statement of Part~I stands. \PI{Remark~9.63}(c) finds the pattern $\so3\to\so2$ with the clock as the
Higgs vector and notes that a clock field with a covariant kinetic term is missing; under (A1) and (A3) the clock vector acts on the spinors
as $\sigma(n\cdot e)$, a vector of $\mathfrak g_1$ that commutes with $\lint$, and it can carry the covariant kinetic term of
\cref{eq:Lcomp}, which is assumed.

The vacuum $T$ of the Dirac phase of \PI{Thm~7.12} is the pole of \cref{prop:higgs1}. The larger kernel orbit of Part~I includes an
$\Omega$ phase. The real-triplet action \cref{eq:coupled} has the vacuum manifold $S^2$ before gauge identification; its two tangent
directions are absorbed by the massive pair, and the remaining scalar fluctuation is radial and massive for $\lambda,Z>0$. The variation of
$T$ along the rotation generated by $\Omega$ is proportional to $T\Omega$, in the other triple of $\mathcal M_1$, along which
\cref{eq:coupled} has no scalar coordinate. Realising the additional phase dynamically would require an enlarged scalar sector.
